\section{DALI Source Patch System}
\label{sec:patch-system}

\subsection{Motivation}

The F1 Race example depends on three core DALI source files located in
\texttt{DALI/src/}:

\begin{itemize}[leftmargin=2em]
  \item \texttt{active\_server\_wi.pl}: starts the LINDA Linda server
  \item \texttt{active\_user\_wi.pl}: starts the interactive user agent
  \item \texttt{active\_dali\_wi.pl}: the main agent runner (\textasciitilde2200 lines)
\end{itemize}

\noindent
These files belong to the upstream DALI framework and are shared across all
examples. Modifying them directly in the repository would couple the framework
to the F1 example, making future upstream merges error-prone.  At the same
time, the F1 example requires a number of fixes that are not yet part of the
upstream code\footnote{I am thinkink of doing a pull request to improve the source code}:

\begin{enumerate}[leftmargin=2em]
  \item \textbf{Dynamic port selection}: the server must try ports from 3010 to 3019
        and write the bound address to \texttt{server.txt} so that every agent
        connects to the correct port, avoiding \texttt{SPIO\_E\_NET\_ADDRINUSE}
        on restart.
  \item \textbf{Address propagation to agents}: \texttt{active\_dali\_wi.pl}
        was hardcoded to \lstinline|'localhost':3010|; it must now read
        \texttt{server.txt} at startup (\texttt{SPIO\_E\_NET\_CONNREFUSED} fix).
  \item \textbf{User agent entry point}: \texttt{active\_user\_wi.pl}
        must read \texttt{server.txt} and connect only when explicitly invoked
        via \lstinline|--goal utente.|, with no top-level directive that could
        block the process on load.
\end{enumerate}

\noindent
The \emph{patch system} solves this tension: the upstream files are patched
\emph{automatically at runtime}, every time the MAS is started, with no
permanent changes stored in the repository.

% ─────────────────────────────────────────────────────────────────────────────
\subsection{Structure}

\begin{lstlisting}[style=tree, numbers=none, basicstyle=\ttfamily\footnotesize]
f1_race/
  patch/
    apply_dali_wi.sh        <- main patch script
    active_server_wi.pl     <- patched copy of the server file
    active_user_wi.pl       <- patched copy of the user-agent file
\end{lstlisting}

\noindent
\texttt{active\_dali\_wi.pl} is \emph{not} copied into \texttt{patch/}
because it is $\approx$ 2200 lines long; instead a short Python snippet
performs a surgical in-place edit.

% ─────────────────────────────────────────────────────────────────────────────
\subsection{Idempotency: the \texttt{[F1\_PATCH]} Marker}

Every patched file contains the following comment on line 7:

\begin{lstlisting}[language=Prolog, numbers=none]
% [F1_PATCH] read server address from server.txt (dynamic port support)
\end{lstlisting}

\noindent
Before applying any change, \texttt{apply\_dali\_wi.sh} checks for this
marker with \texttt{grep}. If it is found, the file is skipped entirely; if
not, the patch is applied and the marker is inserted. This makes every
invocation of the script safe to repeat without double-patching.

\subsection{Reverting the Patches}

To revert all three files to their original upstream state, remove the
\texttt{[F1\_PATCH]} marker line and restore the original code. The fastest
way is via \texttt{git}:

\begin{lstlisting}[language=bash, numbers=none]
git checkout -- src/active_server_wi.pl \
               src/active_user_wi.pl   \
               src/active_dali_wi.pl
\end{lstlisting}

\noindent
The next run of \texttt{startmas.sh} will re-apply the patches automatically.
